% Migrated from the early umbrella Paper-I draft during the five-paper refactor.
% Canonical ownership is now Paper III: Global Knowledge Portrait.
\section{Knowledge absorption and the global portrait}

ORION is intentionally knowledge-absorbent. External disciplines are not extensions that automatically redefine the core; they are projections that may expose new facets, representations, mechanisms, assumptions, contradictions, or search directions.

For an external contribution $c$, absorption returns a typed relation to the current portrait, including equivalent view, complementary facet, refinement, new context coordinate, new representation, new mechanism, contradiction, exposed assumption, exposed search-universe gap, exposed method failure, or unresolved status.

The important transition is therefore not merely
\[
K_t\rightarrow K_{t+1},
\]
but potentially
\[
(K_t,W_t,M_t)\rightarrow(K_{t+1},W_{t+1},M_t),
\]
when new knowledge changes what ORION knows to search for. Changes to $M_t$ remain separately governed.

The global portrait is not a vote or textual average. It should preserve recoverable source projections, context, provenance, incompatibilities and unresolved coordinates while supporting a derived ORION-native representation when that representation explains the compatible local views more economically.
